\section{Discussion}
\subsection{What the simulation evidence establishes}
RQ1 is answered at the implementation-consistency level. The native modal bodies participate in the same dynamics as the robot, the payload mass is not duplicated across lateral coordinates, and archived traces expose forces, actions, and energy diagnostics. The contact correction permits represented cases that previously encountered a guarded near-minimum ambiguity, while rejecting unsupported cases explicitly. Together these features make the simulation inspectable. They do not identify the actual contact compliance or free-surface response of a physical robot carrying water.

RQ2 is partly answered. Full-horizon V4 cases demonstrate simultaneous tracking and modal-domain admissibility on particular mild surfaces. The retained M failure shows why aggregate RMSE alone is insufficient to certify the original task. Teacher experiments further demonstrate that reducing a short forecast energy objective can conflict with globally anchored tracking. A finite action grid can exhaust its performance-feasible options while the plant remains physically admissible. Relaxing performance through an explicitly labeled supervisor is consequently different from silently weakening a safety guard.

The teacher's window-dependent energy sign also has a physical-model interpretation. A lightly damped internal oscillator can exchange energy with commanded base motion; its instantaneous stored energy need not decay monotonically. Lower mean energy over a longer window is not evidence that every prefix improves. Fair reporting should retain the complete energy and tracking histories, use a common observation horizon, and distinguish energy integrals from damping loss or motor work. A future controller objective may need several of these quantities, but this synthesis does not retune the historical reward.

\subsection{Why coverage did not establish learning superiority}
RQ3 remains conditional. Exact-likelihood training and valid gradients remove some implementation explanations for poor learning, but correctness does not imply an informative optimization problem. A useful action held over a local diagnostic horizon may need temporal persistence or state-dependent coordination. Y samples independently at the policy period and preserves the original discount and advantage-estimation parameters. It therefore changes spatial probability coverage while leaving the temporal mechanism untested.

The matched results also reveal a small objective trade-off. Under one initialization, Mixture provides better tracking than Standard at slightly greater modal energy; under the other, both metrics are worse. Selecting the one policy that passes a V4-relative gate does not reverse the direct matched comparison. That selected policy is evidence of a possible local improvement, not an estimator of the average benefit of the exploration family. The short benefit also fails to remain jointly favorable at the conditional three-second aggregate gate.

Several hypotheses remain compatible with these observations: insufficient temporal credit, an unfavorable tracking--energy trade-off in the fixed reward, variance effects independent of mixture geometry, and limited training-state coverage. None is established as the unique cause. A controlled variance-matched or temporally persistent comparison could distinguish them, but would be a new experiment requiring its own protocol. Simply increasing training steps after a mixed development result would not resolve causal ambiguity.

\subsection{Information and comparison fairness}
The complete controller's privileged information must remain visible. V4 uses true simulated modal state; the teacher additionally uses native prediction. The Y actor also receives modal features. Their advantage cannot be presented as robot-only intelligence or a comparison of equally observable policies. Earlier BC and PPO checkpoints use other task and observation contracts and cannot be relabeled as trained figure-eight controllers because their action dimension matches.

Likewise, a censored contact episode should be reported as an environment-domain limitation, while a completed trajectory that misses a registered tracking window is a task failure. These categories answer different scientific questions. Preserving unfavorable categories may narrow the headline result, but it improves the value of the archive for subsequent research and prevents an apparently comprehensive comparison from resting on incompatible denominators.
